% =====================================================================
%  CMPT 981 Project - Milestone (NeurIPS 2022 format)
%
%  Build (from this folder):  latexmk milestone.tex
%  latexmkrc puts the PDF in ../docs/ and keeps build files here.
%  References live in ref.bib (shared with the proposal).
%
%  Paper assignments (one per author):
%    Sec. 2.1  Mroueh (2025)                -> Azaan Rehman
%    Sec. 2.2  Suk & Duan (2025)            -> Ankush Singh
%    Sec. 2.3  Thrampoulidis et al. (2026)  -> Simon Doan
%  Lines starting with % are guidance only; delete them as you write.
% =====================================================================

\documentclass{article}

% to compile a camera-ready version, add the [final] option, e.g.:
\usepackage[final]{neurips_2022}

\usepackage[utf8]{inputenc} % allow utf-8 input
\usepackage[T1]{fontenc}    % use 8-bit T1 fonts
\usepackage{hyperref}       % hyperlinks
\usepackage{url}            % simple URL typesetting
\usepackage{booktabs}       % professional-quality tables
\usepackage{amsfonts}       % blackboard math symbols
\usepackage{nicefrac}       % compact symbols for 1/2, etc.
\usepackage{microtype}      % microtypography
\usepackage{xcolor}         % colors

% --- Added for this project -------------------------------------------
\usepackage{amsmath, amssymb, amsthm}
\usepackage{mathtools}
\usepackage{bm}
\usepackage{graphicx}
\usepackage{subcaption}
\usepackage{enumitem}
\usepackage[capitalize,noabbrev]{cleveref} % must load after hyperref

% Theorem environments
\newtheorem{theorem}{Theorem}
\newtheorem{proposition}[theorem]{Proposition}
\newtheorem{lemma}[theorem]{Lemma}
\theoremstyle{definition}
\newtheorem{definition}[theorem]{Definition}

% Handy macros
\newcommand{\R}{\mathbb{R}}
\newcommand{\E}{\mathbb{E}}
\newcommand{\1}{\mathbf{1}}
\newcommand{\KL}{\mathrm{KL}}
\newcommand{\pref}{p_{\mathrm{ref}}}
\newcommand{\piref}{\pi_{\mathrm{ref}}}
\DeclareMathOperator{\logit}{logit}
\DeclareMathOperator{\Bin}{Bin}
\DeclareMathOperator*{\argmax}{arg\,max}
\newcommand{\norm}[1]{\left\lVert #1 \right\rVert}

% Inline TODO notes; comment out the first line and uncomment the
% second to hide them all before submitting.
\newcommand{\todo}[1]{\textcolor{red}{[TODO: #1]}}
% \newcommand{\todo}[1]{}


\title{GRPO unraveled: A Small-Scale Empirical Study of RLVR\\
  CMPT 981 Project Milestone}


\author{%
  Ankush Singh\\
  Simon Fraser University\\
  \texttt{asa715@sfu.ca} \\
  \And
  Simon Doan\\
  Simon Fraser University\\
  \texttt{vtd@sfu.ca} \\
  \And
  Azaan Rehman\\
  Simon Fraser University\\
  \texttt{ara199@sfu.ca} \\
}


\begin{document}


\maketitle

\begin{abstract}
  \todo{Abstract. Write last: the project goal in one sentence, the three
  papers reviewed, and the progress made since the proposal.}
\end{abstract}

% ---------------------------------------------------------------------
\section{Introduction}
\label{sec:intro}

% - One short paragraph recapping the project goal from the proposal.
% - Why these three papers: each one gives a prediction that one of our
%   core experiments tests (Mroueh -> E2, Suk & Duan -> E3,
%   Thrampoulidis et al. -> E4).
% - What the rest of the milestone covers.
\todo{Introduction}

% ---------------------------------------------------------------------
\section{Paper Reviews}
\label{sec:reviews}

% Shared setup used by all three reviews, so each review can refer to it
% instead of repeating it: prompt x, policy pi_theta, binary reward
% r in {0, 1}, success probability p, group size G, advantages A_i,
% KL strength beta.
\todo{Shared notation (a few lines)}

% Each review uses the same four parts so the three read consistently.
% Keep each review to a similar length.

% ---------------------------------------------------------------------
\subsection{\texorpdfstring{Success amplification~\citep{mroueh2025rlvr}}{Success amplification (Mroueh)}}
\label{sec:mroueh}
% Reviewer: Azaan Rehman

\paragraph{Setting and assumptions.}
% Binary rewards, reverse-KL penalty, no clipping; which quantities are
% assumed exact (expectations, the inner KL-regularized problem).
\todo{Setting and assumptions}

\paragraph{Main results.}
% GRPO's effective loss as a contrastive loss between correct and
% incorrect samples; the recurrence p_n = h(p_{n-1}); fixed points
% p* > p_ref (success amplification); the condition on beta for local
% convergence; mean-only vs. mean+std normalization.
\todo{Main results}

\paragraph{Connection to our project.}
% What E2 tests: finite G vs. the paper's G -> infinity, SGD vs. the
% exact iteration, reverse vs. forward KL. State the prediction we
% compare against.
\todo{Connection to our project}

\paragraph{Limitations and open questions.}
\todo{Limitations and open questions}

% ---------------------------------------------------------------------
\subsection{\texorpdfstring{Gradient gap and step-size thresholds~\citep{suk2025gradientgap}}{Gradient gap and step-size thresholds (Suk and Duan)}}
\label{sec:suk}
% Reviewer: Ankush Singh

\paragraph{Setting and assumptions.}
\todo{Setting and assumptions}

\paragraph{Main results.}
% Definition of the gradient gap; the sufficient step-size condition and
% the lower-bound constructions; scaling with response length T and
% success rate (1 - p); the logistic growth relation; what the paper says
% about GRPO's length normalization.
\todo{Main results}

\paragraph{Connection to our project.}
% What E3 tests: the learning-rate sweep, the 1/L scaling, and the
% mixed-length case where prompts share one learning rate.
\todo{Connection to our project}

\paragraph{Limitations and open questions.}
\todo{Limitations and open questions}

% ---------------------------------------------------------------------
\subsection{\texorpdfstring{Pass@$k$ advantage shaping~\citep{thrampoulidis2026passk}}{Pass@k advantage shaping (Thrampoulidis et al.)}}
\label{sec:thrampoulidis}
% Reviewer: Simon Doan

\paragraph{Setting and assumptions.}
\citet{thrampoulidis2026passk} study how advantage shaping changes the
objective optimized by RLVR methods. They consider binary rewards, where
a correct response receives 1 and 0 otherwise. Under their setting, for
each prompt $x$, the model independently samples a training group of $G$
responses from its current policy. If its success probability is
$p_\theta(x)$, the probability of obtaining at least one correct response
in $k$ independent attempts is
\[
  P_k(x) = 1 - (1 - p_\theta(x))^k.
\]
Here, $G$ is the number of responses during the training phase, while $k$
is the number of attempts considered by the pass@$k$ evaluation metric.
Their analysis assumes one policy update per freshly sampled training
group, with clipping inactive and KL regularization omitted.

\paragraph{Main results.}
The gradient of the $P_k(x)$, pass@$k$ objective, is
\[
  \nabla_\theta P_k(x)
  = k(1 - p_\theta(x))^{k-1}\nabla_\theta p_\theta(x).
\]
For fixed $k > 1$, a high success probability $p_\theta(x)$ makes
$1 - p_\theta(x)$ really small. So pass@$k$ optimization places less
weight on prompts already solved reliably, and more weight to prompts
with lower success probabilities.

To estimate the pass@$k$ gradient from sampled responses, the authors
construct $\mathrm{RLOO}_k$, an unbiased estimator that uses the other
responses in each group to calculate leave-one-out weights. Applying
standard deviation normalization produces their biased $\mathrm{GRPO}_k$
variant. More broadly, the paper connects advantage shaping to surrogate
reward maximization, explaining how existing shaping rules can be
interpreted through the objective they optimize.

\paragraph{Connection to our project.}
For our experiment, we are aiming to compare GRPO with biased and
unbiased pass@$k$ estimators on a simple synthetic task that has one
correct answer. To do this, we will vary the training group size $G$,
while holding the total sampling budget fixed. We will then measure how
each trained model performs on different pass@$k$ versions. Our goal here
is to determine whether larger training groups improve performance under
the same sampling budget. Moreover, under these settings, we can also
compute population gradients directly and use it to compare them with
average sampled gradients during training.

\paragraph{Limitations and open questions.}
\citet{thrampoulidis2026passk} note that their surrogate reward
interpretation relies on an infinite group limit with $k$ being fixed.
But for practical perspectives and also in our experiments, we use
finite groups instead, where noisy success rate estimates can cause
updates to differ from the theoretical prediction. This raises a
question about how closely training with small groups follows the
population predictions.

% ---------------------------------------------------------------------
\iffalse
TODO: If we wanna add these parts in:
\section{Progress}
\label{sec:progress}

% What is implemented so far (e.g. the bandit playground in main.py),
% any preliminary results, and how they compare with the predictions
% above. Replace the placeholder figure with real plots.
\todo{Progress since the proposal}

\begin{figure}[t]
  \centering
  % Replace the box with: \includegraphics[width=0.6\linewidth]{figures/<name>.pdf}
  \fbox{\parbox[c][4cm][c]{0.6\linewidth}{\centering Figure placeholder}}
  \caption{\todo{Caption}}
  \label{fig:prelim}
\end{figure}

% ---------------------------------------------------------------------
\section{Updated Plan}
\label{sec:plan}

% Changes to the experiments since the proposal (E1--E6), anything
% dropped or added, and the timeline to the final report.
\todo{Remaining work and timeline}

% ---------------------------------------------------------------------
\section{Contributions}
\label{sec:contributions}

\begin{itemize}
  \item \textbf{Ankush Singh:} review of \citet{suk2025gradientgap}. \todo{other work}
  \item \textbf{Simon Doan:} review of \citet{thrampoulidis2026passk}. \todo{other work}
  \item \textbf{Azaan Rehman:} review of \citet{mroueh2025rlvr}. \todo{other work}
\end{itemize}

\fi

% ---------------------------------------------------------------------
\bibliographystyle{plainnat}
\bibliography{ref}

\end{document}
